\section{Experimental Setup}
\label{sec:experiments}

\subsection{Experimental Question}

\textbf{Does per-sample last-layer gradient cosine similarity with the within-batch mean gradient achieve higher ROC-AUC than per-sample loss as a predictor of spurious-minority group membership on Waterbirds and CelebA?}

Secondary questions: (1) How does the signal evolve across epochs? (2) Is failure consistent across datasets with different minority prevalences? (3) What does the gap magnitude suggest about the failure mechanism?

\subsection{Datasets}

\textbf{Waterbirds} \citep{Wah2011Caltech,Sagawa2020Distributionally}: Synthetic spurious correlation (bird species $\times$ background, 95\% strength). Train: 4,795 samples, ${\sim}82\%$ majority. Minority for evaluation: landbird on water; waterbird on land (groups \{1, 2\} in group\_DRO encoding).

\textbf{CelebA} \citep{Liu2015Deep}: Blond hair prediction; spurious attribute: gender. Train: 16,000 samples (stratified subsample, ${\sim}10\%$ of full 162K dataset, seed 42). Minority: blond male (${\sim}0.8\%$ of training data, group ID 3).

CelebA complements Waterbirds with much more extreme minority prevalence --- key for testing whether prevalence modulates the contamination effect.

\subsection{Baselines}

\begin{itemize}
  \item \textbf{Gradient Alignment Score ($a_i$):} Negated within-batch cosine similarity (Section~\ref{sec:methodology})
  \item \textbf{Per-Sample Loss ($\ell_i$):} Cross-entropy loss; signal used by JTT and LfF
\end{itemize}

Identical model states; single variable is signal type.
We evaluate signal discriminability (ROC-AUC) rather than downstream worst-group accuracy; the latter is only relevant if the signal exists.

\subsection{Evaluation}

\textbf{ROC-AUC} (sklearn): threshold-free discriminability measure.
AUC $< 0.5$ indicates inverted signal.
Binary labels: $y_i = \mathbb{1}[\text{group\_id}_i \in \text{minority}]$.
